% !TeX root = ../sections/03_solutions.tex

\subsection{1_2}
\begin{solution}
	周期 \(2\pi\) の関数 \(f(x)=x\) のフーリエ級数を求める。
	
	一般に，周期 \(2\pi\) の関数のフーリエ級数は
	\[
	f(x)\sim \frac{a_0}{2}
	+\sum_{n=1}^{\infty}
	\left(a_n\cos nx+b_n\sin nx\right)
	\]
	であり，フーリエ係数は
	\[
	a_0=\frac{1}{\pi}\int_{-\pi}^{\pi}f(x)\,dx,
	\]
	\[
	a_n=\frac{1}{\pi}\int_{-\pi}^{\pi}f(x)\cos nx\,dx,
	\]
	\[
	b_n=\frac{1}{\pi}\int_{-\pi}^{\pi}f(x)\sin nx\,dx
	\]
	で与えられる。
	
	ここで \(f(x)=x\) は奇関数である。
	また \(\cos nx\) は偶関数，\(\sin nx\) は奇関数である。
	
	したがって，
	\[
	x\cos nx
	\]
	は奇関数であり，
	\[
	x\sin nx
	\]
	は偶関数である。
	
	よって，
	\[
	a_0=0,\qquad a_n=0
	\]
	である。
	
	したがって，あとは \(b_n\) を求めればよい。
	\[
	b_n
	=
	\frac{1}{\pi}\int_{-\pi}^{\pi}x\sin nx\,dx
	\]
	である。
	
	被積分関数 \(x\sin nx\) は偶関数なので，
	\[
	b_n
	=
	\frac{2}{\pi}\int_0^\pi x\sin nx\,dx
	\]
	となる。
	
	ここで部分積分を用いる。
	\[
	u=x,\qquad dv=\sin nx\,dx
	\]
	とおくと，
	\[
	du=dx,\qquad v=-\frac{\cos nx}{n}
	\]
	である。
	
	したがって，
	\[
	\int_0^\pi x\sin nx\,dx
	=
	\left[-\frac{x\cos nx}{n}\right]_0^\pi
	+
	\frac{1}{n}\int_0^\pi \cos nx\,dx.
	\]
	
	第2項は
	\[
	\int_0^\pi \cos nx\,dx
	=
	\left[\frac{\sin nx}{n}\right]_0^\pi
	=0
	\]
	である。
	
	よって，
	\[
	\int_0^\pi x\sin nx\,dx
	=
	-\frac{\pi\cos n\pi}{n}.
	\]
	
	ここで
	\[
	\cos n\pi=(-1)^n
	\]
	だから，
	\[
	\int_0^\pi x\sin nx\,dx
	=
	-\frac{\pi(-1)^n}{n}.
	\]
	
	したがって，
	\[
	b_n
	=
	\frac{2}{\pi}
	\left(
	-\frac{\pi(-1)^n}{n}
	\right)
	=
	\frac{2(-1)^{n+1}}{n}.
	\]
	
	以上より，
	\[
	x\sim
	\sum_{n=1}^{\infty}
	\frac{2(-1)^{n+1}}{n}\sin nx
	\]
	である。
	
	したがって，\(5x\) の項まで書くと，
	\[
	\boxed{
		x\sim
		2\left(
		\sin x
		-\frac{1}{2}\sin 2x
		+\frac{1}{3}\sin 3x
		-\frac{1}{4}\sin 4x
		+\frac{1}{5}\sin 5x
		+\cdots
		\right)
	}
	\]
	である。
\end{solution}